\chapter{Summary and Conclusions}
\label{ch:conclusions}

\section{Conclusions}

This work has addressed the design of a high-level quantum programming language and the implementation of its compiler, starting from a central observation laid out in \autoref{sec:contexto}, that no existing language combines a static linear type system with output to a standard, interoperable format supported by multiple hardware vendors. The two design objectives that followed from that observation have guided the entire development of Qalc.

The first objective, giving the language a linear type system without the performance penalty introduced by solutions embedded in a host language such as Guppy, has been met by combining Hindley-Milner type inference with the single-consumption discipline described in \autoref{ch:analisis-semantico}. By being implemented as a purpose-built compiler rather than a library embedded in a dynamically typed language, Qalc verifies at compile time both accidental cloning and unreleased qubit leaks, without requiring explicit type annotations on most expressions of the program. Type-discriminated function overloading, described in \autoref{sec:overloading}, complements this type system by allowing simple polymorphic operations to be expressed without resorting to type classes.

The second objective, generating code towards a standard, interoperable format, has been met with the code generation phase described in \autoref{ch:backend}, which translates the flattened program into OpenQASM~3. The design of the emission mechanism, based on the \lstinline|Transpilable| trait parameterized over the backend type, was deliberately conceived so that adding a further output format would not require modifying the rest of the compiler, in line with the modularity and extensibility goal that runs through the whole architecture described in \autoref{ch:arquitectura}.

The combined result of both objectives is a functional prototype in which the compiler frontend, from lexical analysis through type inference and recursion verification, is fully implemented. The backend phases (the lazy flattener, the linear consumption check and the code generation, described in \autoref{ch:backend}) remain at an experimental stage, sufficient to demonstrate the complete operation of the pipeline over programs that respect the linear discipline described in \autoref{sec:disciplina_lineal_uso}, but they do not yet cover every possible combination of the language.

As anticipated in \autoref{sec:alcance}, this work does not address automatic uncomputation. This is a deliberate limitation rather than an unmet objective, since that problem requires a static analysis of the topological graph of the quantum state that exceeds the scope of a prototype of this kind and constitutes, as already argued, a research line of its own.

\section{Future work}

The modular design of the compiler, in which each phase communicates with the next through the \lstinline|CompilerStep| trait described in \autoref{ch:arquitectura}, makes it possible to identify future lines of work that can be tackled incrementally without redesigning the existing compiler.

\begin{itemize}
  \item \textit{Automatic uncomputation.} Extend the semantic analysis to automatically synthesize the uncomputation operator for recursive quantum data structures, as discussed in \autoref{sec:alcance}, combining static scope analysis with a richer representation of the entanglement graph.

  \item \textit{Maturing the backend phases.} Complete the coverage of the lazy flattener over nested combinations of higher-order combinators and monadic blocks, and extend the test coverage of the linear consumption check to the nested-thunk cases described in \autoref{sec:linear-checker}, until it reaches the same level of maturity as the compiler frontend.

  \item \textit{Expanding the standard library of types and operations.} The delivered prototype offers an illustrative catalog of built-ins (the basic quantum gates, a small set of arithmetic operations and the functions used in this document's examples), sufficient to demonstrate the language design but not for productive use. Completing this standard library relies on the same \lstinline|BuiltinDef| and \lstinline|Transpilable| mechanism already described in \autoref{ch:backend}, without requiring structural changes to the compiler.

  \item \textit{Additional backends.} Exploit the separation between the \lstinline|Transpilable| trait and its concrete implementation for OpenQASM~3, described in \autoref{sec:codegen}, to add support for further output formats, such as Quil, validating in practice the extensibility that motivated that design.

  \item \textit{More expressive function overloading.} Generalize the overloading discriminated by the type of the first parameter, described in \autoref{sec:overloading}, towards resolution mechanisms that consider several parameters or that approach type classes, as long as type inference remains decidable.

  \item \textit{Empirical evaluation on quantum hardware.} Run the circuits generated by Qalc on simulators and on real physical hardware from different providers compatible with OpenQASM~3, to measure the depth and fidelity of the resulting circuits against equivalent implementations in Qiskit or Cirq.
\end{itemize}
